%=============================================================================
% FairTox -- DLE-AI-202 final project, defence deck (Track 2)
%
% Build:  cd presentation && latexmk -pdf presentation.tex
%
% RULES THIS DECK FOLLOWS -- do not undo them:
%
% 1. Minimum text. A slide carries a figure or a table and at most a few short
%    lines. The evidence is the object on the slide; the speaker is the prose.
% 2. Figures come from results/<run>/figures/ wherever one exists -- they are
%    already screen-sized with a title baked in. Only the two cross-run figures
%    (s_families, s_dose) are rendered by scripts/make_slide_figures.py, because
%    no run produces them. Tables come from scripts/make_slide_tables.py, which
%    imports the paper's formatters, so no number here is typed by hand.
% 3. Each speaker's block is CONTIGUOUS: nobody returns later in the deck.
%    Shared frames (title, roadmap, limitations, conclusion) carry no name --
%    the team answers those together.
%
%   1-2    shared             title, roadmap
%   3-6    Rustamova, Nigar   the problem, the three families, the corpus, the
%                             classical baseline
%   7-10   Mirzayeva, Laman   the three axes, the guards, what the errors are,
%                             two coded examples
%   11-14  Samadov, Nijat     the model, the compute and the noise floor,
%                             counterfactual swapping and its null result
%   15-19  Aliyev, Orkhan     the mechanism, the loss, the headline, every
%                             family, the three controls
%   20-21  shared             limitations, conclusion
%   22-23  backup             the alpha sweep, the per-subgroup audit
%
% Every frame was checked as a rendered page: nothing may run under the
% footline, which a LaTeX compiler will not always tell you about.
%=============================================================================
\documentclass[aspectratio=169,11pt,xcolor={table}]{beamer}

\usepackage[T1]{fontenc}
\usepackage[utf8]{inputenc}
\usepackage{helvet}
\usepackage{booktabs}
\usepackage{graphicx}
\usepackage{amsmath}
\usepackage{array}

% The deck reads the run figures where they already live.
\graphicspath{{figures/}{../results/main/figures/}}

%-- palette -------------------------------------------------------------------
\definecolor{fairblue}{HTML}{3B5BA5}
\definecolor{fairgreen}{HTML}{2F8F5B}
\definecolor{fairwarn}{HTML}{C0562A}
\definecolor{fairgrey}{HTML}{6B7280}
\definecolor{fairink}{HTML}{222222}
\definecolor{fairwash}{HTML}{F3F4F6}

\usetheme{default}
\usecolortheme{default}
\setbeamertemplate{navigation symbols}{}
\setbeamercolor{normal text}{fg=fairink}
\setbeamercolor{structure}{fg=fairblue}
\setbeamercolor{frametitle}{fg=fairblue}
\setbeamerfont{frametitle}{size=\large,series=\bfseries}
\setbeamerfont{framesubtitle}{size=\small}
\setbeamercolor{framesubtitle}{fg=fairgrey}
\setbeamertemplate{blocks}[rounded][shadow=false]
\setbeamercolor{block title}{bg=fairblue!14,fg=fairblue}
\setbeamercolor{block body}{bg=fairblue!5,fg=fairink}
\setbeamercolor{block title alerted}{bg=fairwarn!14,fg=fairwarn}
\setbeamercolor{block body alerted}{bg=fairwarn!5,fg=fairink}
\setbeamercolor{block title example}{bg=fairgreen!14,fg=fairgreen}
\setbeamercolor{block body example}{bg=fairgreen!5,fg=fairink}
\setbeamertemplate{itemize item}{\color{fairblue}\rule[0.25ex]{0.7ex}{0.7ex}}
\setbeamertemplate{itemize subitem}{\color{fairgrey}\rule[0.25ex]{0.5ex}{0.5ex}}

\setbeamertemplate{frametitle}{%
  \vspace{1.1ex}%
  \begin{beamercolorbox}[wd=\textwidth]{frametitle}%
    \usebeamerfont{frametitle}\insertframetitle%
    \ifx\insertframesubtitle\@empty\else
      \\[0.2ex]{\usebeamerfont{framesubtitle}\usebeamercolor[fg]{framesubtitle}%
      \insertframesubtitle}%
    \fi
  \end{beamercolorbox}%
  \vspace{0.3ex}%
  {\color{fairblue!35}\rule{\textwidth}{0.6pt}}%
  \vspace{-0.6ex}%
}

% Without \leavevmode this template typesets at \textwidth, reports an overfull
% hbox on every frame, and pushes the body down until it runs under the footline.
\newcommand{\thespeaker}{}
\setbeamercolor{footline}{fg=fairgrey}
\setbeamertemplate{footline}{%
  \leavevmode%
  \hbox{%
  \begin{beamercolorbox}[wd=\paperwidth,ht=2.4ex,dp=1.4ex,leftskip=2.2ex,%
    rightskip=2.2ex]{footline}%
    \scriptsize\insertsection\hfill\thespeaker\hfill\insertframenumber%
  \end{beamercolorbox}}%
  \vskip0pt%
}

\newcommand{\repourl}{https://github.com/nigarrustamova/fairtox}
\newcommand{\repotag}{v1.0-final}

\newcommand{\takeaway}[1]{%
  \vspace{0.6ex}%
  \begin{beamercolorbox}[wd=\textwidth,sep=1.0ex,rounded=true]{block body example}%
    \small\textbf{#1}%
  \end{beamercolorbox}}

% Every figure gets a width AND a height cap: the height is what keeps a caption
% from being pushed under the footline, and the compiler does not warn when that
% happens.
\newcommand{\figfit}[3][0.88]{%
  \begin{center}%
    \includegraphics[width=#1\textwidth,height=#2\textheight,keepaspectratio]{#3}%
  \end{center}}

\newcommand{\head}[1]{{\color{fairblue}\textbf{#1}}}
\newcommand{\hint}[1]{{\footnotesize\color{fairgrey}#1}}

\title{Parity, Utility and Safety Are Not Independent}
\date{DLE-AI-202 \quad Cohort I, 2026}

\begin{document}

%=============================================================================
% 1 -- title (shared)
%=============================================================================
\section{FairTox}
\begin{frame}[plain]
  \vspace{3.5ex}
  \begin{center}
    {\Large\bfseries\color{fairblue} Parity, Utility and Safety Are Not
      Independent\par}
    \vspace{1.0ex}
    {\normalsize A Controlled Comparison of Three Bias-Mitigation Families\par}
    {\normalsize for Toxicity Classification\par}
    \vspace{1.8ex}
    {\color{fairblue!35}\rule{0.55\textwidth}{0.6pt}\par}
    \vspace{2.0ex}
    {\small Aliyev, Orkhan \qquad Mirzayeva, Laman \qquad Rustamova, Nigar
      \qquad Samadov, Nijat\par}
    \vspace{1.4ex}
    {\small AI Academy \,·\, National Artificial Intelligence Centre \,·\,
      Baku\par}
    \vspace{0.5ex}
    \hint{DLE-AI-202, Cohort I 2026 \,·\, Track 2: applied / domain research}
    \vspace{2.4ex}

    \hint{\texttt{\repourl} \,·\, tag \texttt{\repotag}}
  \end{center}
\end{frame}

%=============================================================================
% 2 -- roadmap (shared)
%=============================================================================
\begin{frame}{The talk in four parts}
  \vspace{2.0ex}
  \small
  \begin{center}
  \begin{tabular}{@{}c@{\hspace{2.4em}}l@{\hspace{3.2em}}l@{}}
    \head{1} & The problem, and the corpus & Rustamova, Nigar \\[1.1ex]
    \head{2} & How we measure, and what the errors are & Mirzayeva, Laman \\[1.1ex]
    \head{3} & The model, and the obvious fix & Samadov, Nijat \\[1.1ex]
    \head{4} & The mechanism, and what actually works & Aliyev, Orkhan \\
  \end{tabular}
  \end{center}

  \vspace{2.6ex}
  \begin{block}{The question}
    One corpus, one architecture, one recipe: how do mitigations at the three
    intervention points compare on parity --- and what does each cost in utility
    and in safety?
  \end{block}

  \vspace{0.8ex}
  \centerline{\hint{$38$ runs \,·\, $115$ models \,·\, $34.4$ GPU-hours \,·\,
    every number here is generated from committed artefacts}}
\end{frame}

%=============================================================================
% 3-6 -- Nigar
%=============================================================================
\renewcommand{\thespeaker}{Rustamova, Nigar}
\section{1 · The problem and the corpus}

\begin{frame}{Two errors --- and they land on different people}
  \vspace{2.4ex}
  \begin{columns}[t,onlytextwidth]
    \begin{column}{0.46\textwidth}
      \begin{alertblock}{False negative}
        Abuse gets through.\\[0.5ex]
        The person it targets is \textbf{unprotected}.
      \end{alertblock}
    \end{column}
    \begin{column}{0.46\textwidth}
      \begin{block}{False positive}
        A benign comment is removed.\\[0.5ex]
        The person who wrote it is \textbf{silenced}.
      \end{block}
    \end{column}
  \end{columns}

  \vspace{3.0ex}
  \centerline{\small Identity terms co-occur with abuse in scraped text --- so
    the \emph{token} becomes cheap evidence.}

  \takeaway{Then moderation is a tax on the communities it was deployed to
    protect.}
\end{frame}

%=============================================================================
\begin{frame}{Three places to intervene --- rarely compared}
  \vspace{2.0ex}
  \begin{columns}[t,onlytextwidth]
    \begin{column}{0.30\textwidth}
      \begin{block}{\textsc{pre} --- the data}
        \small Rewrite the text so the term stops predicting the label.\\[0.6ex]
        \hint{Dixon 2018 · Park 2018}
      \end{block}
    \end{column}
    \begin{column}{0.30\textwidth}
      \begin{block}{\textsc{in} --- the loss}
        \small Change the objective: price the errors that matter.\\[0.6ex]
        \hint{Zhang 2018 · Sagawa 2020}
      \end{block}
    \end{column}
    \begin{column}{0.30\textwidth}
      \begin{block}{\textsc{post} --- the threshold}
        \small Leave the model; move the boundary per group.\\[0.6ex]
        \hint{Hardt 2016}
      \end{block}
    \end{column}
  \end{columns}

  \vspace{3.0ex}
  \centerline{\small Each is normally evaluated \textbf{alone}, against its own
    baseline, on its own metric.}

  \takeaway{We run all three: five mitigations, three controls, eight seeds on
    the headline pair.}
\end{frame}

%=============================================================================
\begin{frame}{The corpus: Jigsaw Unintended Bias}
  \vspace{2.2ex}
  \begin{center}
    \small
    \begin{tabular}{lr}
\toprule
Quantity & Value \\
\midrule
Comments in the release & 1,999,516 \\
Identity-annotated \emph{(our population)} & 447,998 \; ($22.4\%$) \\
Toxic rate & $0.1134$ \\
\addlinespace
Split: train / val / test & 335,998 / 44,800 / 67,200 \\
\addlinespace
Subgroups configured & 24 \\
Clearing the FPR / FNR floor & 14 / 9 \\
\bottomrule
\end{tabular}

  \end{center}
  \vspace{2.2ex}
  \small
  \begin{itemize}
    \item Target is the annotator \emph{fraction} --- binarised at $0.5$, audit
      re-run at $\tau = 0.3 \ldots 0.7$.
    \item Identity-annotated rows only: unrated $\neq$ identity-free. Licence
      CC0 1.0.
  \end{itemize}
\end{frame}

%=============================================================================
\begin{frame}{A baseline worth beating}
  \vspace{2.4ex}
  \begin{center}
    \small
    \begin{tabular}{lrrrr}
\toprule
Model & macro-F1 & ROC-AUC & FPR & FNR \\
\midrule
TF-IDF $+$ logistic regression (balanced) & $0.7460$ & $0.9096$ & $0.1179$ & $0.2352$ \\
DistilBERT, fine-tuned & $0.8003$ & $0.9387$ & $0.0289$ & $0.4211$ \\
\bottomrule
\end{tabular}

  \end{center}
  \vspace{2.4ex}
  \centerline{\small Balanced class weights on purpose --- unweighted logistic
    regression reaches FNR $0.97$.}

  \takeaway{Very different operating points: we quote both error rates, never
    macro-F1 alone.}
\end{frame}

%=============================================================================
% 7-10 -- Laman
%=============================================================================
\renewcommand{\thespeaker}{Mirzayeva, Laman}
\section{2 · How we measure}

\begin{frame}{Three axes --- never combined into one score}
  \vspace{2.2ex}
  \begin{columns}[t,onlytextwidth]
    \begin{column}{0.30\textwidth}
      \begin{block}{Utility}
        \small macro-F1, ROC-AUC, PR-AUC\\[0.6ex]
        \hint{is it a good classifier?}
      \end{block}
    \end{column}
    \begin{column}{0.30\textwidth}
      \begin{block}{Parity}
        \small FPR gap, macro FPR, EOD\\[0.6ex]
        \hint{who is wrongly censored?}
      \end{block}
    \end{column}
    \begin{column}{0.30\textwidth}
      \begin{block}{Safety}
        \small subgroup FNR\\[0.6ex]
        \hint{whose abuse is missed?}
      \end{block}
    \end{column}
  \end{columns}

  \vspace{3.0ex}
  \centerline{$\mathrm{FPR}_k = \Pr(\hat p \geq \tau \mid y = 0,\, k)$
    \qquad gap $= \max_k \mathrm{FPR}_k - \min_k \mathrm{FPR}_k$}

  \takeaway{A single ``is it better?'' number would hide the trade-off we exist
    to measure.}
\end{frame}

%=============================================================================
\begin{frame}{Five guards, all fixed before we saw a result}
  \vspace{2.4ex}
  \small
  \begin{itemize}\itemsep 1.5ex
    \item \head{Separate support floors} ($N \geq 100$ each) ---
      \textbf{14 subgroups for parity, 9 for safety}.
    \item \head{Wilson intervals} --- a bootstrap returns $[0,0]$ at zero false
      positives, whatever the sample size.
    \item \head{Collapse check} --- a model that flags nothing has
      \textbf{perfect parity}. None of ours collapsed.
    \item \head{One split per seed}, fingerprinted; duplicates measured:
      $\leq 1.11\%$ of test rows, effect $\leq 0.0032$.
    \item \head{Outcomes A / B / C / D pre-registered} in a config file before
      any model was trained.
  \end{itemize}
\end{frame}

%=============================================================================
\begin{frame}{What the model actually gets wrong}
  \vspace{1.6ex}
  \begin{center}
    \small
    \begin{tabular}{lrrr}
\toprule
What the comment actually is & $n$ & \% & Label arguable \\
\midrule
hostile, but not aimed at the group & 21 & $35.0$ & 11 \\
genuinely derogatory towards the group & 16 & $26.7$ & 16 \\
argues \emph{against} bigotry & 12 & $20.0$ & 0 \\
neutral mention, no hostility & 7 & $11.7$ & 0 \\
quotes hostility in order to criticise it & 4 & $6.7$ & 0 \\
\bottomrule
\end{tabular}

  \end{center}
  \vspace{0.6ex}
  \centerline{\hint{a uniform random sample of $60$ identity-bearing false
    positives, read and coded by hand}}
  \vspace{1.4ex}
  \begin{columns}[t,onlytextwidth]
    \begin{column}{0.47\textwidth}
      \begin{alertblock}{$45\%$ arguable}
        \small a reasonable annotator could have called them toxic
      \end{alertblock}
    \end{column}
    \begin{column}{0.47\textwidth}
      \begin{block}{$38\%$ protected speech}
        \small counter-speech, quotation, neutral mention
      \end{block}
    \end{column}
  \end{columns}
\end{frame}

%=============================================================================
\begin{frame}{Two of the sixty}
  \vspace{3.0ex}
  \begin{block}{\normalsize ``He was mentally-ill.''}
    \small model score $0.984$ \,·\, dataset label: non-toxic \,·\,
    \textbf{neutral mention}
  \end{block}

  \vspace{2.0ex}
  % The comment continues; the ellipsis says so rather than presenting the first
  % sentence as the whole of it.
  \begin{block}{\normalsize ``Yes, they condemn Islamic terrorists, who are no
    different than Nazis and White Supremacists\,\ldots''}
    \small model score $0.789$ \,·\, dataset label: non-toxic \,·\,
    \textbf{counter-speech}
  \end{block}

  \takeaway{The worst error a moderation system can make: silencing the speech
    the policy exists to protect.}
\end{frame}

%=============================================================================
% 11-14 -- Nijat
%=============================================================================
\renewcommand{\thespeaker}{Samadov, Nijat}
\section{3 · The model, and the obvious fix}

\begin{frame}{The model}
  \vspace{2.4ex}
  \begin{columns}[t,onlytextwidth]
    \begin{column}{0.46\textwidth}
      \begin{block}{Architecture}
        \small DistilBERT, pretrained\\
        \texttt{[CLS]} $\rightarrow$ dropout $\rightarrow$
        \texttt{Linear}$(768,1)$\\[0.6ex]
        \hint{DistilRoBERTa for the backbone arm --- no code change}
      \end{block}
    \end{column}
    \begin{column}{0.48\textwidth}
      \begin{block}{Recipe}
        \small AdamW · lr $3\times10^{-5}$ · batch $64$ · $3$ epochs\\
        max length $128$ · mixed precision\\[0.6ex]
        \hint{best-epoch selection, early stopping --- identical in both arms}
      \end{block}
    \end{column}
  \end{columns}

  \vspace{2.6ex}
  \takeaway{No \texttt{AutoModelForSequenceClassification}: it hides the loss,
    and every intervention here lives in the loss.}
\end{frame}

%=============================================================================
\begin{frame}{115 models --- and how much of a difference is noise}
  \vspace{2.4ex}
  \begin{columns}[t,onlytextwidth]
    \begin{column}{0.46\textwidth}
      \begin{block}{The budget we used}
        \small One MIG $3\text{g}.20\text{gb}$ slice of an A100\\
        $0.064$--$0.071$ s/step · peak $3.99$ GB · $34.4$ GPU-h\\[0.6ex]
        \hint{spare memory went to more seeds, not a bigger batch}
      \end{block}
    \end{column}
    \begin{column}{0.48\textwidth}
      \begin{block}{The noise floor}
        \small The same baseline trained \textbf{five} times\\
        lands in \textbf{two} states, not five\\[0.6ex]
        \hint{$46$ of $67{,}200$ test decisions differ --- $0.068\%$}
      \end{block}
    \end{column}
  \end{columns}

  \vspace{2.6ex}
  \takeaway{Floor $0.0000$ on the FPR gap and $0.0017$ on the worst-moving
    metric --- about $78\times$ smaller than the effect we report.}
\end{frame}

%=============================================================================
\begin{frame}{The obvious fix: swap the identity term}
  \vspace{3.0ex}
  \begin{center}
    \normalsize
    \emph{``I am a proud \textbf{Muslim} woman''} $\;\longrightarrow\;$
    \emph{``I am a proud \textbf{Christian} woman''}
  \end{center}
  \vspace{3.0ex}
  \small
  \begin{itemize}\itemsep 1.4ex
    \item Swapped \textbf{in place}, never appended --- otherwise the step count
      changes too.
    \item Two doses: $p = 0.5$ and $p = 1.0$, five seeds each.
    \item Assumes toxicity is invariant under the swap --- true for abuse, false
      for facts.
  \end{itemize}
\end{frame}

%=============================================================================
\begin{frame}{It does not work --- and the way it fails is the result}
  \figfit[0.84]{0.505}{s_dose}
  \vspace{-0.6ex}
  \centerline{\hint{five seeds per dose; the slope is the shortcut regression of
    the next section}}
  \takeaway{At full dose the shortcut flattens $30.9\%$ and the gap still does
    not close: the token was never the whole cause.}
\end{frame}

%=============================================================================
% 15-19 -- Orkhan
%=============================================================================
\renewcommand{\thespeaker}{Aliyev, Orkhan}
\section{4 · The mechanism, and what works}

% The figure carries its own title, so the frame title must not repeat it.
\begin{frame}{So where does the gap come from?}
  \vspace{1.0ex}
  \begin{columns}[c,onlytextwidth]
    \begin{column}{0.58\textwidth}
      \includegraphics[width=\linewidth,height=0.70\textheight,keepaspectratio]%
        {fig10_mechanism}
    \end{column}
    \begin{column}{0.40\textwidth}
      \small
      A subgroup's \textbf{toxic rate} predicts the model's \textbf{FPR} on that
      subgroup almost perfectly.\\[1.4ex]
      $r = 0.885$--$0.925$ in \emph{every} configuration we ran.\\[1.4ex]
      \begin{beamercolorbox}[wd=\linewidth,sep=1.0ex,rounded=true]%
        {block body example}%
        \small\textbf{Not prejudice --- a conditional probability learned from
        the corpus and applied as a prior.}%
      \end{beamercolorbox}
    \end{column}
  \end{columns}
\end{frame}

%=============================================================================
\begin{frame}{The fix: one \emph{asymmetric} weight}
  \vspace{2.6ex}
  \begin{columns}[t,onlytextwidth]
    \begin{column}{0.48\textwidth}
      \centering
      \small
      \renewcommand{\arraystretch}{1.6}
      \begin{tabular}{l|c|c}
        \multicolumn{1}{c}{} & \multicolumn{1}{c}{\textbf{non-toxic}}
          & \multicolumn{1}{c}{\textbf{toxic}} \\ \cline{2-3}
        \textbf{background} & \cellcolor{fairwash}$1$ & \cellcolor{fairwash}$1$ \\ \cline{2-3}
        \textbf{identity-bearing} & \cellcolor{fairgreen!30}$\boldsymbol{\alpha}$
          & \cellcolor{fairwash}$1$ \\ \cline{2-3}
      \end{tabular}

      \vspace{2.2ex}
      $\displaystyle \mathcal{L} = \frac{\sum_i w_i \ell_i}{\sum_i w_i}$
      \qquad $\alpha = 4$
    \end{column}
    \begin{column}{0.46\textwidth}
      \vspace{-1.4ex}
      \small
      \begin{itemize}\itemsep 1.3ex
        \item One \textsc{train} function; the arms differ in \textbf{one
          argument}.
        \item Split and training fingerprints must match --- otherwise the
          comparison is \textbf{refused}.
        \item $\alpha$ fixed in advance, swept in the ablation.
      \end{itemize}
    \end{column}
  \end{columns}
\end{frame}

%=============================================================================
\begin{frame}{The headline, over eight seeds}
  \vspace{1.8ex}
  \begin{center}
    \small
    \setlength{\tabcolsep}{3pt}
    \begin{tabular}{llrrrlc}
\toprule
Axis & Metric & Baseline & Mitigated & $\Delta$ & 95\% CI & Seeds \\
\midrule
Parity & subgroup FPR gap & $0.1072$ & $0.0552$ & $-0.0520$ & $[-0.0742,\,-0.0297]$ & 8/8 \\
Utility & macro-F1 & $0.8031$ & $0.7903$ & $-0.0128$ & $[-0.0153,\,-0.0104]$ & 8/8 \\
Safety & subgroup FNR & $0.4382$ & $0.5520$ & $+0.1138$ & $[+0.0824,\,+0.1452]$ & 8/8 \\
\bottomrule
\end{tabular}

  \end{center}
  \vspace{1.8ex}
  \small
  \begin{itemize}\itemsep 1.0ex
    \item $8/8$ seeds agree in sign · sign test $p = 0.0078$ · pre-registered
      outcome \textbf{C}.
    \item The shortcut slope flattens $44.6\%$ --- and the fit stays just as
      tight.
    \item Threshold-free metrics agree: BPSN $+0.0114$, BNSP $-0.0169$.
  \end{itemize}
  \takeaway{The price: $2.19$ points of missed abuse per point of FPR gap
    closed.}
\end{frame}

%=============================================================================
\begin{frame}{The three intervention points differ in \emph{sign}}
  \figfit[0.86]{0.66}{s_families}
  \vspace{-0.4ex}
  \centerline{\hint{filled marker = the paired $95\%$ interval excludes zero;
    three-seed families carry a direction only}}
\end{frame}

%=============================================================================
\begin{frame}{Three controls, one conclusion}
  \vspace{1.0ex}
  \begin{center}
    \footnotesize
    \begin{tabular}{lcrrrr}
\toprule
Control: what it weights & Seeds & \multicolumn{2}{c}{FPR gap} & \multicolumn{2}{c}{subgroup FNR} \\
\cmidrule(lr){3-4}\cmidrule(lr){5-6}
 & & before & after & before & after \\
\midrule
inverse class frequency (\texttt{class}) & 3 & $0.1152$ & $0.3005$ & $0.4183$ & $0.2039$ \\
identity-toxic cell alone & 1 & $0.0767$ & $0.1374$ & $0.4744$ & $0.3374$ \\
Group DRO, weights \emph{learned} & 3 & $0.1153$ & $0.3056$ & $0.4186$ & $0.1877$ \\
\bottomrule
\end{tabular}

  \end{center}
  \vspace{1.2ex}
  \begin{columns}[t,onlytextwidth]
    \begin{column}{0.46\textwidth}
      \centering
      \footnotesize
      \setlength{\tabcolsep}{4pt}
      \begin{tabular}{lrr}
\toprule
Training cell & Ours & Learned \\
\midrule
background, non-toxic & $0.143$ & $0.000$ \\
background, toxic & $0.143$ & $0.150$ \\
identity, non-toxic & $0.571$ & $0.440$ \\
identity, toxic & $0.143$ & $0.410$ \\
\bottomrule
\end{tabular}

    \end{column}
    \begin{column}{0.50\textwidth}
      \vspace{-0.6ex}
      \begin{alertblock}{The claim}
        \small The gain needs an \textbf{asymmetric} weight on the identity
        \emph{non-toxic} cell alone. Raise the \emph{toxic} cell too --- by hand
        or learned --- and the frontier runs backwards.
      \end{alertblock}
    \end{column}
  \end{columns}
\end{frame}

%=============================================================================
% 20-21 -- shared close
%=============================================================================
\renewcommand{\thespeaker}{}
\section{Limitations and conclusion}

\begin{frame}{What we cannot claim}
  \vspace{2.0ex}
  \small
  \begin{itemize}\itemsep 1.4ex
    \item Absolute parity improved; \texttt{white}/\texttt{christian} relative
      parity did \textbf{not} ($3.32\times \rightarrow 3.83\times$).
    \item One subgroup got worse: \texttt{atheist} $+50\%$ from a small base.
    \item $45\%$ of the coded false positives are \textbf{annotator
      disagreement} --- $n = 60$, one coder.
    \item Support is uneven: $14$ of $24$ subgroups for parity, $9$ for safety.
    \item Only two of eight configurations reach eight seeds.
    \item One corpus, one language, coarse binary identity categories.
  \end{itemize}
\end{frame}

%=============================================================================
\begin{frame}{Conclusion}
  \vspace{2.0ex}
  \small
  \begin{itemize}\itemsep 1.5ex
    \item The disparity is a \textbf{property of the training distribution}, not
      a bug in one model.
    \item It can be reduced --- $-0.0520$ over eight seeds --- but never for
      free: $2.19$ points of missed abuse per point.
    \item The three families differ in \textbf{sign}, not merely in strength.
    \item Three controls say why: the gain needs an \emph{asymmetric} weight, so
      even a learned worst-group objective cannot find it.
  \end{itemize}

  \vspace{2.2ex}
  \begin{beamercolorbox}[wd=\textwidth,sep=1.1ex,rounded=true]%
    {block body example}%
    \centering\small\textbf{Thank you --- questions?}\\[0.3ex]
    \hint{\texttt{\repourl} \,·\, tag \texttt{\repotag}}%
  \end{beamercolorbox}
\end{frame}

%=============================================================================
% 22-23 -- backup
%=============================================================================
\appendix
\section{Backup}
\renewcommand{\thespeaker}{backup}

\begin{frame}{How strong should the weighting be?}
  \figfit[0.78]{0.60}{s_alpha}
  \vspace{-0.4ex}
  \centerline{\hint{three seeds; parity flattens after the pre-registered
    $\alpha = 4$ while the safety cost keeps climbing}}
\end{frame}

\begin{frame}{Who the intervention moves}
  \figfit[0.88]{0.66}{fig05_fpr_by_subgroup}
  \vspace{-0.4ex}
  \centerline{\hint{$12$ of the $14$ improve; \texttt{asian} is unchanged and
    \texttt{atheist} --- the smallest group --- gets worse}}
\end{frame}

\end{document}
